\section{Key Management During Transition}


\subsection{Four Key Types}

During Phases 1 and 2, validators maintain four keys simultaneously:

\begin{center}
\begin{tabular}{lllll}
\toprule
\textbf{Key} & \textbf{Algorithm} & \textbf{Storage} & \textbf{Rotation} & \textbf{HSM Support} \\
\midrule
EC & secp256k1 & Keystore file & Manual & Yes (PKCS\#11) \\
BLS & BLS12-381 & Keystore file & Per-epoch capable & Yes (custom) \\
ML-DSA & ML-DSA-65 & Keystore file & Per-epoch capable & Planned (FIPS 140-3) \\
Ringtail & Ring-LWE & DKG shares & Per-epoch (automatic) & N/A (distributed) \\
\bottomrule
\end{tabular}
\end{center}

\subsection{Key Hierarchy}

\begin{enumerate}[leftmargin=2em]
\item The \textbf{ML-DSA key} is the validator's long-term identity. It certifies all other keys.
\item The \textbf{BLS key} is certified by the ML-DSA key at registration. Rotation requires a new ML-DSA certificate.
\item The \textbf{EC key} derives the validator's EVM address. It is bound to the ML-DSA key via a registration transaction.
\item The \textbf{Ringtail key} is generated via DKG and rotated automatically each epoch. It is ephemeral.
\end{enumerate}

\subsection{Validator Upgrade Path}

Existing validators upgrade through three steps:

\begin{enumerate}[leftmargin=2em]
\item \textbf{Generate ML-DSA key pair.} Offline, on validator hardware.
\item \textbf{Register ML-DSA public key.} P-Chain transaction signed by the existing EC key, binding the ML-DSA key to the validator's identity.
\item \textbf{Upgrade node software.} The new node software signs blocks with both BLS and ML-DSA keys. No downtime required---the upgrade is a rolling restart.
\end{enumerate}

The entire process takes $<$5 minutes per validator and requires no coordination between validators.

